% TOP_PROBLEM: {"id": 30007435, "title": "Schrodinger convergence at the higher-dimensional endpoint", "category_id": 8, "category": {"id": 8, "name": "analysis", "display_name": "Analysis", "description": "Limits, continuity, calculus, and function theory.", "slug": "analysis", "order_index": 8, "created_at": "2026-07-31T15:26:25.671Z"}, "status": "open"}
% FIRST_POSTED: 2026-10-06
% ATTEMPT_STATUS: solved
% OPENAI_CLAIM: {"schema_version":1,"source_repo":"https://github.com/openai/math","source_commit":"adc7f1241b42e322a6451854ab7e4b4c146bf78a","release_date":"2026-10-06","match":"full","resolution":"proved","independently_reviewed":false,"second_pass":"agreed","families":[{"family":"080","title":"The exact Sobolev endpoint for Schrödinger convergence","subject":"Real and complex analysis","manuscripts":[{"title":"Endpoint pointwise convergence for the Schrodinger equation in higher dimensions","pdf_path":"preprints/Endpoint-pointwise-convergence-for-the-Schrodinger-equation-in-higher-dimensions-September-24-2026/paper.pdf","readme_path":"preprints/Endpoint-pointwise-convergence-for-the-Schrodinger-equation-in-higher-dimensions-September-24-2026/README.md","date":"2026-09-24","theorem_ref":"preprints/Endpoint-pointwise-convergence-for-the-Schrodinger-equation-in-higher-dimensions-September-24-2026/build/sections/introduction.tex :: m:main"}],"lean":null,"reasoning_trace":null}]}
\documentclass[11pt]{article}
\usepackage[margin=25mm]{geometry}
\usepackage{fontspec}
\setmainfont{Latin Modern Roman}
\usepackage{amsmath,amssymb,mathtools}
\usepackage{unicode-math}
\setmathfont{Latin Modern Math}
\usepackage{enumitem}
\usepackage{xurl,hyperref,bookmark}
\hypersetup{colorlinks=true,urlcolor=blue}
\setcounter{secnumdepth}{0}
\setlength{\parindent}{0pt}
\setlength{\parskip}{5pt}
\setlength{\emergencystretch}{3em}
\begin{document}
\section*{Schrodinger convergence at the higher-dimensional endpoint}

\subsection{OpenAI's claim}
\textit{Higher-dimensional Schrödinger endpoints}: For every $n\ge3$, the claimed critical Sobolev order for almost-everywhere free Schrödinger convergence is $n/(2(n+1))$. The regularized evolution is asserted to have a common spatial representative throughout the positive time interval.

\subsection{Outline of the argument}
\textit{Higher-dimensional Schrödinger endpoints}: Transverse packet fields are controlled by a frame count with polynomial dependence on transversality, combined with amplitude and occupancy hierarchies. Induction produces a short-time maximal estimate at precisely the critical frequency factor. A single time tree is used across frequencies for square summation, after which Gaussian approximation removes exceptional sets uniformly in time.

\subsection{Scope relative to this problem}
\textit{Higher-dimensional Schrödinger endpoints}: The primary statement says: For every n>=3 and every f in H\textasciicircum{}(n/(2(n+1)))(Rn), Gaussian frequency regularizations of Schrodinger evolution have a representative continuous down to zero at almost every point, on a common full-measure set. The target retains these quantified objects and hypotheses, so its positive assertion is covered.

With the hypotheses and conventions above, the claimed result addresses the whole mathematical assertion on this page.

\subsection{Formal verification}
\textit{Higher-dimensional Schrödinger endpoints}: The pinned catalogue lists no Lean formalization for this result.

This site has not run the Lean development. These links describe OpenAI’s selected formal statements and should be read with their stated scope.

\subsection{Status caveat}
OpenAI reports that an internal model produced the manuscripts. Their arguments have not been independently verified as proofs by this site. Any portion outside the documented formal scope may contain mistakes, and the catalogue label reports the announced mathematical claim.

\subsection{Sources}
\begin{itemize}
\item OpenAI, \emph{Endpoint pointwise convergence for the Schrodinger equation in higher dimensions}, 2026-09-24; citation key \texttt{OAI:Endpoint-pointwise-convergence-for-the-Schrodinger-equation-in-higher-dimensions-September-24-2026}. \href{https://github.com/openai/math/blob/adc7f1241b42e322a6451854ab7e4b4c146bf78a/preprints/Endpoint-pointwise-convergence-for-the-Schrodinger-equation-in-higher-dimensions-September-24-2026/paper.pdf}{PDF (GitHub)}; \href{https://github.com/openai/math/raw/adc7f1241b42e322a6451854ab7e4b4c146bf78a/preprints/Endpoint-pointwise-convergence-for-the-Schrodinger-equation-in-higher-dimensions-September-24-2026/paper.pdf}{PDF (direct)}; \href{https://github.com/openai/math/blob/adc7f1241b42e322a6451854ab7e4b4c146bf78a/preprints/Endpoint-pointwise-convergence-for-the-Schrodinger-equation-in-higher-dimensions-September-24-2026/README.md}{Manuscript page and citation}; \href{https://github.com/openai/math/blob/main/preprints/Endpoint-pointwise-convergence-for-the-Schrodinger-equation-in-higher-dimensions-September-24-2026/paper.pdf}{Latest version}.
\item \href{https://github.com/openai/math/blob/adc7f1241b42e322a6451854ab7e4b4c146bf78a/CONTENTS.md}{OpenAI catalogue entry}.
\end{itemize}
\end{document}
